\documentclass{umk-matematika}

\begin{document}

\maketitlepage
    {Практическая работа №1}
    {Радианная мера углов}
    {}

\section{Фундамент (1--4 балла)}

\textbf{Задача 1 (1 балл).} Переведите 30$^\circ$ в радианы.

\textbf{Задача 2 (2 балла).} Переведите 90$^\circ$ в радианы.

\textbf{Задача 3 (3 балла).} Переведите $\dfrac{\pi}{4}$ радиан в градусы.

\textbf{Задача 4 (4 балла).} Переведите 180$^\circ$ в радианы.

\section{Стены (5--7 баллов)}

\textbf{Задача 5 (5 баллов).} Переведите 150$^\circ$ в радианы.

\textbf{Задача 6 (6 баллов).} Переведите $\dfrac{2\pi}{3}$ радиан в градусы.

\textbf{Задача 7 (7 баллов).} Заполните пропуск: 270$^\circ$ = \_\_\_ рад.

\section{Кровля (8--10 баллов)}

\textbf{Задача 8 (8 баллов).} Переведите 315$^\circ$ в радианы. Ответ дайте в виде дроби.

\textbf{Задача 9 (9 баллов).} Угол наклона кровли 40$^\circ$. Выразите в радианах (округлите до сотых, $\pi \approx 3{,}14$).

\textbf{Задача 10 (10 баллов).} Кран повернулся на $\dfrac{5\pi}{6}$ радиан. Сколько это в градусах?

\newpage
\section*{Ответы}

\begin{enumerate}
  \item $\dfrac{\pi}{6}$
  \item $\dfrac{\pi}{2}$
  \item 45$^\circ$.
  \item $\pi$
  \item $\dfrac{5\pi}{6}$
  \item 120$^\circ$.
  \item $\dfrac{3\pi}{2}$
  \item $\dfrac{7\pi}{4}$
  \item ≈ 0,70 рад.
  \item 150$^\circ$.
\end{enumerate}

\end{document}